(1) 合并指数得到
\[
E(x)=E_0\exp[-\tfrac12(w^{-2}-ib)x^2].
\]
因 \(\Re(w^{-2}-ib)=w^{-2}>0\)，复高斯积分收敛，完成平方给
\[
\widetilde E(q)
=E_0\sqrt{\frac{2\pi}{w^{-2}-ib}}
\exp\left[-\frac{q^2}{2(w^{-2}-ib)}\right].
\]
平方根选择从 \(b=0\) 连续延伸的分支，使 \(\widetilde E(0)=E_0\sqrt{2\pi}\,w\) 当 \(b=0\)；其他支仅给整体符号但需和物理相位约定一致。

(2) 令 \(s=bw^2\)，由
\[
\Re\frac1{w^{-2}-ib}=\frac{w^2}{1+s^2}
\]
可得归一化远场强度
\[
\frac{I(q)}{I(0)}
=\exp\left[-\frac{w^2q^2}{1+b^2w^4}\right].
\]
因此远场强度降至峰值 \(1/e\) 时的正空间频率半宽为
\[
q_{1/e}=\frac{\sqrt{1+b^2w^4}}{w}.
\]
在窄角近轴区 \(q\simeq k\theta\)，\(\theta_{1/e}\simeq q_{1/e}/k\)。若 \(b=0\)，标准高斯光束角谱宽度为 \(1/w\)，光束越宽，远场衍射越窄。

(3) 输入强度与 \(e^{-x^2/w^2}\) 成正比，故其均方根宽度为 \(\Delta x=w/\sqrt2\)。输出强度与 \(e^{-w^2q^2/(1+s^2)}\) 成正比，因此
\[
\Delta q=\frac{\sqrt{1+s^2}}{\sqrt2\,w},\qquad
\Delta x\,\Delta q=\frac12\sqrt{1+b^2w^4}\ge\frac12.
\]
当且仅当 \(b=0\) 时取最小值 \(1/2\)。二次相位不会改变输入强度宽度，却增加相空间的线性位置—频率关联，使远场角谱变宽；这不与傅里叶不确定性关系矛盾。

(4) 在近轴薄透镜模型中，焦距为 \(f\) 的理想透镜赋予透射场的相位为 \(\exp[-ikx^2/(2f)]\)，因此它相当于 \(b=-k/f\)（若只考虑透镜造成的二次相位）。当 \(|b|w^2\gg1\)，远场角谱半宽趋于 \(q_{1/e}\sim|b|w\)，由相位斜率 \(|\partial_x(bx^2/2)|\sim |b|w\) 可以独立理解。但若 \(q_{1/e}\) 不再远小于 \(k\)，或透镜像差和有限孔径重要，则 \(q\simeq k\theta\) 及一维近轴模型失效。
